\section{Ekstensi Diffie-Hellman}
\label{sec:ekstensi-dh}

Protokol dua pihak yang diuraikan pada Subbagian~\ref{sec:protokol-dh} dapat
diperluas ke beberapa arah. Bagian ini membahas tiga di antaranya: perluasan
kepada lebih dari dua pihak, pemakaian kunci sesaat yang memberi
\textit{forward secrecy}, dan wujud protokol ini pada sistem yang dipakai
sehari-hari.

\subsection{Pertukaran Kunci Tiga Pihak}

Gagasan Diffie--Hellman tidak terbatas pada dua pihak. Untuk Alice, Bob, dan
Carol yang ingin menyepakati satu kunci bersama, pertukaran dilakukan secara
melingkar: setiap pihak memangkatkan nilai yang diterimanya dengan kunci
privatnya sendiri, lalu meneruskannya kepada pihak berikutnya. Setelah satu
putaran penuh, setiap pihak telah memangkatkan nilai awal dengan keempat hal
yang diperlukan, sehingga semuanya sampai pada $K = g^{abc} \bmod p$.

Prosedur lengkapnya terdiri atas dua belas langkah.

\begin{enumerate}
    \item Alice, Bob, dan Carol menyepakati $p$ dan $g$ secara terbuka.
    \item Alice memilih kunci privat $a$, Bob memilih $b$, dan Carol memilih $c$.
          Ketiganya dirahasiakan.
    \item Alice menghitung $g^a \bmod p$ dan mengirimkannya kepada Bob.
    \item Bob menghitung $(g^a)^b \bmod p = g^{ab} \bmod p$ dan mengirimkannya
          kepada Carol.
    \item Carol menghitung $K = (g^{ab})^c \bmod p = g^{abc} \bmod p$.
          Carol kini memegang kuncinya.
    \item Bob menghitung $g^b \bmod p$ dan mengirimkannya kepada Carol.
    \item Carol menghitung $(g^b)^c \bmod p = g^{bc} \bmod p$ dan mengirimkannya
          kepada Alice.
    \item Alice menghitung $K = (g^{bc})^a \bmod p = g^{bca} \bmod p
          = g^{abc} \bmod p$. Alice kini memegang kuncinya.
    \item Carol menghitung $g^c \bmod p$ dan mengirimkannya kepada Alice.
    \item Alice menghitung $(g^c)^a \bmod p = g^{ca} \bmod p$ dan
          mengirimkannya kepada Bob.
    \item Bob menghitung $K = (g^{ca})^b \bmod p = g^{cab} \bmod p
          = g^{abc} \bmod p$. Bob kini memegang kuncinya.
    \item Alice, Bob, dan Carol kini memiliki kunci rahasia yang sama, yaitu
          $K = g^{abc} \bmod p$.
\end{enumerate}

Urutan pada langkah 3 sampai 11 memperlihatkan pola yang berulang tiga kali.
Setiap pengulangan dimulai dengan kunci publik salah satu pihak, dikalikan
berturut-turut oleh dua pihak lainnya, dan berakhir pada nilai $g^{abc}$.
Karena perkalian pangkat bersifat komutatif, ketiga jalur itu bertemu pada
bilangan yang sama, sebagaimana diperlihatkan pada
Gambar~\ref{fig:dh-tiga-pihak}.

% Pertukaran kunci Diffie-Hellman tiga pihak, melingkar dalam tiga putaran.
% Dirujuk dari section-04.tex dengan % Pertukaran kunci Diffie-Hellman tiga pihak, melingkar dalam tiga putaran.
% Dirujuk dari section-04.tex dengan % Pertukaran kunci Diffie-Hellman tiga pihak, melingkar dalam tiga putaran.
% Dirujuk dari section-04.tex dengan \input{figures/dh-tiga-pihak}.
\begin{figure}[htbp]
    \centering
    \resizebox{0.98\linewidth}{!}{%
    \begin{tikzpicture}[
        pihak/.style={draw, rounded corners=1pt, minimum size=1.5cm,
                      align=center, font=\small, fill=blue!6, inner sep=2pt},
        panah/.style={-{Latex[length=2.4mm]}, thick},
        langkah/.style={font=\scriptsize, fill=white, inner sep=2pt,
                        align=center, midway},
    ]
        \node[pihak] (a) at (90:3.1)  {Alice\\$a$};
        \node[pihak] (b) at (210:3.1) {Bob\\$b$};
        \node[pihak] (c) at (330:3.1) {Carol\\$c$};

        % putaran 1: g^a -> Bob -> g^ab -> Carol
        \draw[panah] (a) to[bend right=22]
            node[langkah] {$g^{a}$} (b);
        \draw[panah] (b) to[bend right=22]
            node[langkah] {$g^{ab}$} (c);
        \draw[panah] (c) to[bend right=22]
            node[langkah] {$g^{abc} = K$} (a);

        % putaran 2 dan 3 ditandai di luar lingkaran
        \node[font=\footnotesize, align=left, text width=13.6cm, anchor=north]
            at (0,-4.2)
            {Putaran berikutnya berjalan dengan pola sama: $g^{b} \to g^{bc} \to
             g^{abc}$ untuk Alice, lalu $g^{c} \to g^{ca} \to g^{abc}$ untuk Bob.
             Setiap pihak memangkatkan nilai yang diterimanya dengan kunci
             privatnya sendiri, lalu meneruskannya kepada pihak berikutnya};
    \end{tikzpicture}}
    \caption{Pertukaran kunci tiga pihak: satu putaran melingkar menghasilkan
    $g^{abc}$, dan dua putaran berikutnya dengan pola sama}
    \label{fig:dh-tiga-pihak}
    \sumber{Diolah dari Munir, \textit{IF4020 Kriptografi}, ITB.}
\end{figure}
.
\begin{figure}[htbp]
    \centering
    \resizebox{0.98\linewidth}{!}{%
    \begin{tikzpicture}[
        pihak/.style={draw, rounded corners=1pt, minimum size=1.5cm,
                      align=center, font=\small, fill=blue!6, inner sep=2pt},
        panah/.style={-{Latex[length=2.4mm]}, thick},
        langkah/.style={font=\scriptsize, fill=white, inner sep=2pt,
                        align=center, midway},
    ]
        \node[pihak] (a) at (90:3.1)  {Alice\\$a$};
        \node[pihak] (b) at (210:3.1) {Bob\\$b$};
        \node[pihak] (c) at (330:3.1) {Carol\\$c$};

        % putaran 1: g^a -> Bob -> g^ab -> Carol
        \draw[panah] (a) to[bend right=22]
            node[langkah] {$g^{a}$} (b);
        \draw[panah] (b) to[bend right=22]
            node[langkah] {$g^{ab}$} (c);
        \draw[panah] (c) to[bend right=22]
            node[langkah] {$g^{abc} = K$} (a);

        % putaran 2 dan 3 ditandai di luar lingkaran
        \node[font=\footnotesize, align=left, text width=13.6cm, anchor=north]
            at (0,-4.2)
            {Putaran berikutnya berjalan dengan pola sama: $g^{b} \to g^{bc} \to
             g^{abc}$ untuk Alice, lalu $g^{c} \to g^{ca} \to g^{abc}$ untuk Bob.
             Setiap pihak memangkatkan nilai yang diterimanya dengan kunci
             privatnya sendiri, lalu meneruskannya kepada pihak berikutnya};
    \end{tikzpicture}}
    \caption{Pertukaran kunci tiga pihak: satu putaran melingkar menghasilkan
    $g^{abc}$, dan dua putaran berikutnya dengan pola sama}
    \label{fig:dh-tiga-pihak}
    \sumber{Diolah dari Munir, \textit{IF4020 Kriptografi}, ITB.}
\end{figure}
.
\begin{figure}[htbp]
    \centering
    \resizebox{0.98\linewidth}{!}{%
    \begin{tikzpicture}[
        pihak/.style={draw, rounded corners=1pt, minimum size=1.5cm,
                      align=center, font=\small, fill=blue!6, inner sep=2pt},
        panah/.style={-{Latex[length=2.4mm]}, thick},
        langkah/.style={font=\scriptsize, fill=white, inner sep=2pt,
                        align=center, midway},
    ]
        \node[pihak] (a) at (90:3.1)  {Alice\\$a$};
        \node[pihak] (b) at (210:3.1) {Bob\\$b$};
        \node[pihak] (c) at (330:3.1) {Carol\\$c$};

        % putaran 1: g^a -> Bob -> g^ab -> Carol
        \draw[panah] (a) to[bend right=22]
            node[langkah] {$g^{a}$} (b);
        \draw[panah] (b) to[bend right=22]
            node[langkah] {$g^{ab}$} (c);
        \draw[panah] (c) to[bend right=22]
            node[langkah] {$g^{abc} = K$} (a);

        % putaran 2 dan 3 ditandai di luar lingkaran
        \node[font=\footnotesize, align=left, text width=13.6cm, anchor=north]
            at (0,-4.2)
            {Putaran berikutnya berjalan dengan pola sama: $g^{b} \to g^{bc} \to
             g^{abc}$ untuk Alice, lalu $g^{c} \to g^{ca} \to g^{abc}$ untuk Bob.
             Setiap pihak memangkatkan nilai yang diterimanya dengan kunci
             privatnya sendiri, lalu meneruskannya kepada pihak berikutnya};
    \end{tikzpicture}}
    \caption{Pertukaran kunci tiga pihak: satu putaran melingkar menghasilkan
    $g^{abc}$, dan dua putaran berikutnya dengan pola sama}
    \label{fig:dh-tiga-pihak}
    \sumber{Diolah dari Munir, \textit{IF4020 Kriptografi}, ITB.}
\end{figure}


Sepintas prosedur itu tidak berbiaya banyak: dua belas langkah untuk tiga pihak,
sedangkan untuk dua pihak hanya diperlukan empat. Namun perhatikan bahwa
jumlah langkah itu tumbuh jauh lebih cepat daripada jumlah pesertanya. Untuk
$n$ pihak, setiap pihak harus memangkatkan dan meneruskan nilai yang diterimanya
sebanyak $n-1$ kali, sehingga seluruh protokol memerlukan $n(n-1)$ pengiriman.
Karena itu, protokol melingkar semacam ini hanya praktis untuk kelompok yang
sangat kecil. Untuk kelompok yang besar dipakai pendekatan lain, misalnya
memilih satu pihak sebagai pusat yang membangkitkan kunci sesi lalu mendistribusikannya
kepada pihak lain melalui saluran terenkripsi, atau memakai skema berbasis
struktur pohon yang memangkas jumlah pengiriman.

\subsection{Kunci Sesaat dan \textit{Forward Secrecy}}
\label{subsec:forward-secrecy}

Pada protokol dasar, kunci privat $a$ dan $b$ dapat dipakai berulang kali untuk
banyak sesi. Pemakaian ulang semacam itu, yang disebut Diffie--Hellman statis,
memiliki dua kelemahan. Pertama, bila satu kali saja kunci privat $a$ bocor,
maka seluruh percakapan yang pernah memakainya --- termasuk percakapan yang sudah
lama selesai dan mungkin telah direkam --- dapat dibuka oleh penyerang. Kedua,
pemakaian ulang membuka jalan bagi serangan yang lebih canggih yang mengeksploitasi
hubungan antarsesi.

Praktik modern menghindari kedua kelemahan itu dengan menjadikan kunci privat
sebagai nilai \emph{sesaat} (\textit{ephemeral}): $a$ dan $b$ dibangkitkan baru
untuk setiap sesi dan dibuang begitu sesi berakhir. Protokol dengan sifat ini
disebut \textbf{DHE} (\textit{Diffie--Hellman Ephemeral}). Karena setiap sesi
memakai kunci privat yang berbeda, kompromi satu sesi tidak menolong penyerang
untuk membuka sesi yang lain.

Sifat perlindungan inilah yang disebut \textbf{\textit{forward secrecy}} atau
\textit{perfect forward secrecy}: kebocoran kunci jangka panjang --- misalnya
kunci privat peladen yang dipakai untuk tanda tangan pada sertifikatnya --- tidak
mengakibatkan terbukanya rekaman percakapan yang sudah terjadi. Perhatikan
bahwa perlindungan ini bekerja ke arah belakang, bukan ke depan: forward secrecy
tidak menjamin apa pun bagi sesi yang sedang berjalan pada saat kunci itu bocor,
sebab pada saat itu penyerang tentu dapat menyadap kunci sesi yang sedang
dipakai.

Inilah alasan mengapa TLS 1.3 menghapus seluruh cara pertukaran kunci yang tidak
memberi forward secrecy. Pada versi-versi sebelumnya, pertukaran kunci dapat
dilakukan dengan mengenkripsi kunci sesi memakai kunci publik RSA milik peladen
--- cara yang telah dibahas pada Subbagian~\ref{subsec:hibrida-eksplisit}. Cara
itu memang lebih sederhana, tetapi seluruh catatan percakapan yang pernah
direkam dapat dibuka kembali seketika setelah kunci privat RSA peladen bocor,
tanpa batas waktu. Sejak TLS 1.3, pertukaran kunci wajib memakai
Diffie--Hellman sesaat, dan peran kunci RSA direduksi menjadi menandatangani
nilai pertukaran \citep{rfc8446}.

\subsection{Kurva Eliptik: ECDH}

Cara paling penting untuk memperingan Diffie--Hellman adalah memindahkannya dari
grup perkalian modulo prima ke grup titik pada kurva eliptik. Protokolnya
bernama \textbf{ECDH} (\textit{Elliptic Curve Diffie--Hellman}), dan struktur
langkahnya persis sama dengan yang telah dibahas: kedua pihak memilih kunci
privat, menghitung titik publik dengan mengalikan titik basis, bertukar titik
publik, lalu mengalikan lagi. Yang berubah hanyalah arti "perpangkatan" --- ia
menjadi perkalian skalar pada kurva.

Keuntungannya berupa efisiensi yang besar. ECDH dengan kunci 256 bit memberi
tingkat keamanan yang kira-kira setara dengan Diffie--Hellman modulus 3072 bit,
namun dengan kunci yang dua belas kali lebih pendek dan perhitungan yang jauh
lebih ringan. Karena itu ECDH --- khususnya varian X25519 --- menjadi pilihan
utama pada TLS 1.3 dan pada hampir seluruh protokol modern yang menuntut
pertukaran kunci yang cepat. Dasar matematis kurva eliptik, termasuk cara
memilih titik basis dan menghitung orde grupnya, dibahas pada
Bab~\ref{chap:ecc}.

\subsection{Implementasi Modern}

Protokol Diffie--Hellman dipakai pada hampir setiap koneksi terenkripsi yang
terjadi hari ini, meskipun penggunanya jarang menyadarinya.

\begin{itemize}
    \item \textbf{TLS dan HTTPS.} Setiap kunjungan ke situs dengan awalan
          \texttt{https} diawali dengan pertukaran kunci Diffie--Hellman atau
          ECDH untuk menyepakati kunci sesi. Setelah kunci itu ada, seluruh isi
          halaman dienkripsi dengan AES.
    \item \textbf{SSH.} Protokol untuk masuk ke komputer jarak jauh ini memakai
          pertukaran kunci Diffie--Hellman untuk membentuk saluran terenkripsi,
          lalu mengautentikasi pengguna dengan kata sandi atau kunci publik.
    \item \textbf{IPsec.} Lapisan keamanan pada tingkat jaringan ini memakai
          Diffie--Hellman untuk membentuk asosiasi keamanan antara dua gerbang
          jaringan, sehingga seluruh lalu lintas di antaranya dapat dienkripsi
          tanpa perubahan pada aplikasi.
    \item \textbf{Aplikasi pesan.} Protokol Signal, yang menjadi dasar enkripsi
          pada beberapa aplikasi pesan, menjalankan pertukaran kunci
          Diffie--Hellman secara berlapis dan memperbarui kunci sesi secara
          berkala, sehingga memberi forward secrecy bahkan pada satu percakapan
          yang berlangsung lama.

\end{itemize}

Perhatikan pola yang berulang pada keempat contoh itu: Diffie--Hellman selalu
dipakai untuk memperoleh kunci, dan kunci itu selalu dipakai oleh algoritma
kunci-simetri untuk memproses isi pesan. Inilah bentuk kerja sama antara
kriptografi kunci-publik dan kunci-simetri yang telah diperkenalkan pada
Subbagian~\ref{subsec:hibrida-eksplisit}, kini dengan protokol pertukaran kunci
menggantikan RSA pada tahap penetapan kunci. Protokol pertukaran kunci tidak
pernah menggantikan AES; ia hanya menjadikan kunci AES dapat diperoleh secara
aman tanpa pertemuan tatap muka.
